\documentclass[12pt, a4paper]{article}
\usepackage[T2A]{fontenc}
\usepackage[utf8]{inputenc}
\usepackage[russian]{babel}
\usepackage{amsmath, amssymb}
\usepackage{enumitem}
\usepackage{geometry}
\geometry{top=2cm, bottom=2cm, left=2.5cm, right=2cm}

\begin{document}

\section*{Георгий Игоревич Вольфсон <<Делимость с человеческим лицом>>}

\subsection*{Беседа шестая. Уравнения в целых числах.}

\subsubsection*{Упражнения}

\begin{enumerate}[label=\arabic*., wide=0pt, leftmargin=*]

\item \textbf{Решите уравнение в целых числах: \(3x - 6y = 7\).}

\textbf{Решение:}
Левая часть уравнения: \(3x - 6y = 3(x - 2y)\) делится на 3. Правая часть равна 7, что не делится на 3. Следовательно, уравнение не имеет решений в целых числах.
Ответ: решений нет.

\item \textbf{Решите уравнение в целых числах: \(3x + 2y = 5\).}

\textbf{Решение:}
Заметим, что \((1;1)\) — частное решение. Тогда общее решение имеет вид:
\[
x = 1 - 2k,\quad y = 1 + 3k,\quad k \in \mathbb{Z}.
\]
Проверка: \(3(1-2k) + 2(1+3k) = 3 - 6k + 2 + 6k = 5\).
Ответ: \((1-2k;\; 1+3k)\), где \(k \in \mathbb{Z}\).

\item \textbf{Решите уравнение в целых числах: \(x^2 - y^2 = 2\).}

\textbf{Решение:}
Разложим на множители:
\[
(x-y)(x+y) = 2.
\]
Так как 2 — простое число, возможны 4 системы:
\[
\begin{cases}
x-y = 1,\\
x+y = 2,
\end{cases}
\quad
\begin{cases}
x-y = -1,\\
x+y = -2,
\end{cases}
\quad
\begin{cases}
x-y = 2,\\
x+y = 1,
\end{cases}
\quad
\begin{cases}
x-y = -2,\\
x+y = -1.
\end{cases}
\]
Решая каждую систему, получаем:
\begin{itemize}
\item \(x = 1.5,\; y = 0.5\) — не целые,
\item \(x = -1.5,\; y = -0.5\) — не целые,
\item \(x = 1.5,\; y = -0.5\) — не целые,
\item \(x = -1.5,\; y = 0.5\) — не целые.
\end{itemize}
Ни одна система не даёт целых решений.
Ответ: решений нет.

\item \textbf{Решите уравнение в натуральных числах: \(x^2 + 3xy - 4y^2 = 6\).}

\textbf{Решение:}
Разложим левую часть на множители:
\[
x^2 + 3xy - 4y^2 = (x - y)(x + 4y).
\]
Получаем:
\[
(x - y)(x + 4y) = 6.
\]
Так как \(x, y \in \mathbb{N}\), то \(x+4y > 0\) и \(x+4y \ge 1+4=5\). Перебираем возможные делители числа 6:
\begin{itemize}
\item \(x - y = 1,\; x + 4y = 6\). Складываем: \(2x + 3y = 7\). Из первого \(x = y+1\). Подставляем: \(y+1 + 4y = 6 \Rightarrow 5y = 5 \Rightarrow y = 1,\; x = 2\). Проверка: \(2^2 + 3\cdot2\cdot1 - 4\cdot1^2 = 4 + 6 - 4 = 6\) — подходит.
\item \(x - y = 2,\; x + 4y = 3\). Тогда \(x+4y \ge 1+4=5\), но 3 < 5 — невозможно.
\item \(x - y = 3,\; x + 4y = 2\) — невозможно.
\item \(x - y = 6,\; x + 4y = 1\) — невозможно.
\item Рассмотрим отрицательные делители: \(x - y = -1,\; x + 4y = -6\) — невозможно, так как \(x+4y > 0\).
\end{itemize}
Единственное решение: \((2;1)\).
Ответ: \((2;1)\).

\item \textbf{Решите уравнение в натуральных числах: \(x^2 + 2x + 6 = y^2\).}

\textbf{Решение:}
Оценим:
\[
(x+1)^2 = x^2 + 2x + 1 < x^2 + 2x + 6,
\]
\[
(x+3)^2 = x^2 + 6x + 9 > x^2 + 2x + 6 \quad \text{при } x \ge 1.
\]
Значит,
\[
(x+1)^2 < y^2 < (x+3)^2.
\]
Так как \(x+1, y, x+3\) — натуральные числа, то
\[
x+1 < y < x+3.
\]
Между \(x+1\) и \(x+3\) находится только одно целое число: \(x+2\). Следовательно, \(y = x+2\).
Подставляем в исходное уравнение:
\[
x^2 + 2x + 6 = (x+2)^2 = x^2 + 4x + 4,
\]
\[
2x + 6 = 4x + 4 \Rightarrow 2 = 2x \Rightarrow x = 1.
\]
Тогда \(y = 1 + 2 = 3\).
Ответ: \((1;3)\).

\item \textbf{Решите уравнение в целых числах: \(xy - 2x + y = 7\).}

\textbf{Решение:}
Преобразуем левую часть:
\[
xy - 2x + y = (x+1)(y-2) + 2.
\]
Тогда уравнение принимает вид:
\[
(x+1)(y-2) = 5.
\]
Так как 5 — простое число, возможны 4 варианта:
\begin{enumerate}[label=\arabic*)]
\item \(x+1 = 1,\; y-2 = 5 \Rightarrow x = 0,\; y = 7\);
\item \(x+1 = -1,\; y-2 = -5 \Rightarrow x = -2,\; y = -3\);
\item \(x+1 = 5,\; y-2 = 1 \Rightarrow x = 4,\; y = 3\);
\item \(x+1 = -5,\; y-2 = -1 \Rightarrow x = -6,\; y = 1\).
\end{enumerate}
Ответ: \((0;7),\; (-2;-3),\; (4;3),\; (-6;1)\).

\item \textbf{Решите уравнение в натуральных числах: \(x^2 + y^2 + x^2y^2 = 1969\).}

\textbf{Решение:}
Преобразуем:
\[
x^2 + y^2 + x^2y^2 + 1 = 1970,
\]
\[
(x^2 + 1)(y^2 + 1) = 1970.
\]
Разложим 1970 на простые множители:
\[
1970 = 2 \cdot 5 \cdot 197.
\]
Так как \(x, y \in \mathbb{N}\), то \(x^2 + 1 \ge 2\) и \(y^2 + 1 \ge 2\). Перебираем возможные значения \(x^2 + 1\):
\begin{itemize}
\item \(x^2 + 1 = 2 \Rightarrow x = 1\), тогда \(y^2 + 1 = 985 \Rightarrow y^2 = 984\) — не квадрат.
\item \(x^2 + 1 = 5 \Rightarrow x = 2\), тогда \(y^2 + 1 = 394 \Rightarrow y^2 = 393\) — не квадрат.
\item \(x^2 + 1 = 10 \Rightarrow x = 3\), тогда \(y^2 + 1 = 197 \Rightarrow y^2 = 196 = 14^2 \Rightarrow y = 14\).
\item \(x^2 + 1 = 197 \Rightarrow x^2 = 196 \Rightarrow x = 14\), тогда \(y^2 + 1 = 10 \Rightarrow y = 3\).
\end{itemize}
Также можно рассмотреть \(x^2+1 = 1970\) (тогда \(y=0\), но 0 не натуральное) и другие делители, но они дают не квадраты.
Ответ: \((3;14),\; (14;3)\).

\item \textbf{Решите уравнение в натуральных числах: \(x + \dfrac{1}{y + \frac{1}{z}} = \dfrac{7}{3}\).}

\textbf{Решение:}
Так как \(x \in \mathbb{N}\), а дробь положительна, то \(x < 7/3 < 3\), значит \(x = 1\) или \(x = 2\).

\textbf{Случай 1: \(x = 1\).}
\[
\frac{1}{y + \frac{1}{z}} = \frac{7}{3} - 1 = \frac{4}{3}
\quad\Rightarrow\quad
y + \frac{1}{z} = \frac{3}{4}.
\]
Но \(y \ge 1\), следовательно, \(y + \frac{1}{z} \ge 1 > \frac{3}{4}\). Противоречие.

\textbf{Случай 2: \(x = 2\).}
\[
\frac{1}{y + \frac{1}{z}} = \frac{7}{3} - 2 = \frac{1}{3}
\quad\Rightarrow\quad
y + \frac{1}{z} = 3.
\]
Так как \(y \in \mathbb{N}\), то \(y = 1\) или \(y = 2\).
\begin{itemize}
\item \(y = 1\): \(1 + \frac{1}{z} = 3 \Rightarrow \frac{1}{z} = 2 \Rightarrow z = \frac{1}{2}\) — не натуральное.
\item \(y = 2\): \(2 + \frac{1}{z} = 3 \Rightarrow \frac{1}{z} = 1 \Rightarrow z = 1\) — подходит.
\end{itemize}
Ответ: \((2;2;1)\).

\item \textbf{Решите уравнение в натуральных числах: \((x + y)xy = 120\).}

\textbf{Решение:}
Разложим 120 на множители:
\[
120 = 2^3 \cdot 3 \cdot 5.
\]
Заметим, что \((3;5)\) — решение: \((3+5)\cdot3\cdot5 = 8 \cdot 15 = 120\). В силу симметрии \((5;3)\) также решение.

Докажем, что других решений нет. Пусть \(x \le y\) (в силу симметрии). Тогда \(x^2(x+1) \le (x+y)xy = 120\). Перебираем возможные \(x\):
\begin{itemize}
\item \(x = 1\): \((1+y)\cdot1\cdot y = y(y+1) = 120\). \(10 \cdot 11 = 110 < 120\), \(11 \cdot 12 = 132 > 120\) — решений нет.
\item \(x = 2\): \((2+y)\cdot2\cdot y = 2y(y+2) = 120 \Rightarrow y(y+2) = 60\). \(7 \cdot 9 = 63 > 60\), \(6 \cdot 8 = 48 < 60\) — решений нет.
\item \(x = 3\): \((3+y)\cdot3\cdot y = 3y(y+3) = 120 \Rightarrow y(y+3) = 40\). \(y=5\): \(5 \cdot 8 = 40\) — подходит.
\item \(x = 4\): \((4+y)\cdot4\cdot y = 4y(y+4) = 120 \Rightarrow y(y+4) = 30\). \(y=3\) даёт \(3\cdot7 = 21 < 30\), \(y=4\) даёт \(4\cdot8 = 32 > 30\) — решений нет.
\item \(x = 5\): \(5y(y+5) = 120 \Rightarrow y(y+5) = 24\). \(y=3\) даёт \(3\cdot8 = 24\) — подходит, но это уже \((5;3)\), что симметрично найденному.
\item \(x \ge 6\): \(x^2(x+1) \ge 36 \cdot 7 = 252 > 120\) — невозможно.
\end{itemize}
Таким образом, единственные решения: \((3;5)\) и \((5;3)\).
Ответ: \((3;5),\; (5;3)\).

\item \textbf{Решите уравнение в натуральных числах: \(x^2 + 5xy + y^2 = 7\).}

\textbf{Решение:}
Преобразуем:
\[
x^2 + 2xy + y^2 = 7 - 3xy,
\]
\[
(x+y)^2 = 7 - 3xy.
\]
Так как \(x, y \in \mathbb{N}\), то \(x+y \ge 2\), следовательно, \((x+y)^2 \ge 4\). Также \(xy \ge 1\), значит \(7 - 3xy \le 4\). Таким образом,
\[
4 \le (x+y)^2 = 7 - 3xy \le 4.
\]
Отсюда:
\[
(x+y)^2 = 4 \quad\text{и}\quad 7 - 3xy = 4.
\]
Из первого: \(x+y = 2\) (так как натуральные). Тогда \(x = y = 1\). Проверка: \(1 + 5 + 1 = 7\) — верно.
Ответ: \((1;1)\).

\end{enumerate}

\end{document}